% Options for packages loaded elsewhere
\PassOptionsToPackage{unicode}{hyperref}
\PassOptionsToPackage{hyphens}{url}
%
\documentclass[
  10pt,
]{article}
\usepackage{amsmath,amssymb}
\usepackage{iftex}
\ifPDFTeX
  \usepackage[T1]{fontenc}
  \usepackage[utf8]{inputenc}
  \usepackage{textcomp} % provide euro and other symbols
\else % if luatex or xetex
  \usepackage{unicode-math} % this also loads fontspec
  \defaultfontfeatures{Scale=MatchLowercase}
  \defaultfontfeatures[\rmfamily]{Ligatures=TeX,Scale=1}
\fi
\usepackage{lmodern}
\ifPDFTeX\else
  % xetex/luatex font selection
\fi
% Use upquote if available, for straight quotes in verbatim environments
\IfFileExists{upquote.sty}{\usepackage{upquote}}{}
\IfFileExists{microtype.sty}{% use microtype if available
  \usepackage[]{microtype}
  \UseMicrotypeSet[protrusion]{basicmath} % disable protrusion for tt fonts
}{}
\makeatletter
\@ifundefined{KOMAClassName}{% if non-KOMA class
  \IfFileExists{parskip.sty}{%
    \usepackage{parskip}
  }{% else
    \setlength{\parindent}{0pt}
    \setlength{\parskip}{6pt plus 2pt minus 1pt}}
}{% if KOMA class
  \KOMAoptions{parskip=half}}
\makeatother
\usepackage{xcolor}
\usepackage[margin=1in]{geometry}
\usepackage{longtable,booktabs,array}
\usepackage{calc} % for calculating minipage widths
% Correct order of tables after \paragraph or \subparagraph
\usepackage{etoolbox}
\makeatletter
\patchcmd\longtable{\par}{\if@noskipsec\mbox{}\fi\par}{}{}
\makeatother
% Allow footnotes in longtable head/foot
\IfFileExists{footnotehyper.sty}{\usepackage{footnotehyper}}{\usepackage{footnote}}
\makesavenoteenv{longtable}
\setlength{\emergencystretch}{3em} % prevent overfull lines
\providecommand{\tightlist}{%
  \setlength{\itemsep}{0pt}\setlength{\parskip}{0pt}}
\setcounter{secnumdepth}{-\maxdimen} % remove section numbering
\ifLuaTeX
  \usepackage{selnolig}  % disable illegal ligatures
\fi
\usepackage{bookmark}
\IfFileExists{xurl.sty}{\usepackage{xurl}}{} % add URL line breaks if available
\urlstyle{same}
\hypersetup{
  pdftitle={An actual envelope contact obstruction for every real reference on the six-cone carrier},
  pdfauthor={Unverified P05.2 research draft},
  hidelinks,
  pdfcreator={LaTeX via pandoc}}

\title{An actual envelope contact obstruction for every real reference
on the six-cone carrier}
\author{Unverified P05.2 research draft}
\date{3 October 2026}

\begin{document}
\maketitle

\subsection{Result and exact scope}\label{result-and-exact-scope}

The original P05.2 classification of arbitrary balanced rational
polyhedral pairs remains open. This paper proves an all-reference
negative theorem on one specified two-dimensional carrier. It constructs
an actual PAPC envelope, computes its complete recursive contact
integral, and does not replace a contact calculation by failure of
whole-pair regularity. No supporting theorem is imported from an earlier
note, review, or book. All mathematical lemmas used below are proved
here. This is an unverified draft, with an author audit and finite
diagnostics, without formal verification or priority certification.

Set \(N=\mathbb Z^3\), \(W=N\otimes\mathbb R\), and \[
v_1=(1,0,0),\quad v_2=(0,1,0),\quad v_3=(0,0,1),
\quad v_4=(-1,-1,-1).
\] Let \(X\subset W\) be the union of the six closed cones

\[
\sigma_{ij}=\mathbb R_{\geq0}v_i+\mathbb R_{\geq0}v_j,
\qquad 1\leq i<j\leq4,
\]

with unit top weights. Use the actual saturated tangent lattice

\[
N_\sigma=N\cap\operatorname{span}_{\mathbb R}(\sigma-\sigma)
\]

on every cell, and the inward primitive quotient normals. We prove the
lattice and balancing assertions below.

Here \(\mathrm{PA}_{\mathbb R}\) means continuous functions with real
affine coefficients on a \textbf{finite rational spatial complex}. The
auxiliary class \(\mathrm{PA}_{\mathbb Q}\) requires rational affine
coefficients on such a complex. Neither is the literal BMPS
integral-slope/rational-constant class. PC uses the actual
rational-endpoint polyhedral convex combinations defined below,
including real barycentric coefficients. PAPC additionally satisfies
effectivity against every nonnegative local real balanced cycle on every
ordinary open, in every positive dimension. We establish this property
directly for every PA PC function used in the construction.

\textbf{Theorem (all-real-reference actual contact obstruction; partial
P05.2).} For every

\[
\gamma\in\mathrm{PA}_{\mathbb R}(X)\cap\mathrm{PC}(X)
\]

whose slope is strictly positive on every nonzero recession direction of
every affine reference cell, there are a compactly supported nonnegative

\[
u\in\mathrm{PA}_{\mathbb R,b}(X),
\]

an explicitly constructed \(P\in\mathrm{PAPC}_{\mathbb R}(X,\gamma)\),
rational numbers \(\lambda>0\), \(R>0\), \(S>2R\), a real constant
\(C>0\), and \(\epsilon\in\{1,2\}\), such that

\[
P_\gamma(u)=P,\qquad P=\gamma\quad\text{on }\{\rho\geq S\},
\qquad P=C+\lambda q\quad\text{on }\{\rho\leq2R\},
\]

where

\[
q=|x-y|+|z|,\qquad \rho=|x|+|y|+|z|\quad\text{on }X.
\]

The particular envelope is uniformly decreasingly regularizable, through
\(P+1/n\), in the declared real PAPC domain. Its full finite positive
recursive two-divisor measure satisfies

\[
\int_X\bigl(P-(\gamma+u)\bigr)
\mathrm{MA}_{\mathrm{poly},\mathbb R}(P)
=-8\epsilon\lambda^2<0. \tag{T}
\]

The obstacle can be chosen so that its restriction to \(\sigma_{13}\) is
not a sum of a function of its first cone parameter and a function of
its second cone parameter. In particular, no separability assumption is
imposed or silently used.

Thus none of these references can satisfy the all-obstacle contact
identity, even restricted to obstacles whose envelopes are already PAPC.
This quantifies over all real references in the stated finite
rational-domain class. It does not classify arbitrary carriers,
establish all-envelope regularity on this carrier, or assert a
literal-category or analytic realization result.

\subsection{Primary definitions and dependency
boundary}\label{primary-definitions-and-dependency-boundary}

BMPS denotes A. M. Botero, E. Mazzon and L. Pille-Schneider,
\emph{Monge--Ampère measures on balanced polyhedral spaces},
arXiv:2603.08664v1, submitted 9 March 2026, manuscript dated 10 March
2026. The cached primary PDF and text are snapshotted in
\texttt{sources/}. Its PDF SHA-256 is
\texttt{5cf7e68ee9113f1dc813f9aefa5f9ca72cba02fc4fa688f1a61e4b7f0d1223e7}.
References below are to that version only.

Definitions 2.18--2.20 specify the saturated tangent lattices, inward
primitive normals, and rational weights and cycles. We use the explicit
real extension of weights: a local \(k\)-cycle \(Z\) assigns real
weights to \(k\)-cells, with

\[
\sum_{\sigma\succ\tau} Z(\sigma)n_{\sigma/\tau}
\in N_{\tau,\mathbb R}
\]

at every \((k-1)\)-face. The carrier weight here is integral, namely
one. Allowing real test weights and real affine function coefficients
requires proof, supplied below; it is not inferred from a
rational-weight statement. Definitions 2.22 and Remark 2.37 allow
generalized local fans with lineality and cycles on ordinary open sets.
All our local cycle tests use germs on these opens.

Definition 2.28 is the literal integral-slope/rational-constant function
class. Our rational and real classes are separately declared. On a fixed
rational complex the continuity equations for affine pieces are rational
linear equations; their real solution space has a rational basis.
Clearing denominators makes rational solutions literal. Consequently the
declared real PA class is the real scalar extension on that fixed
spatial structure. Arbitrary real maxima need not belong to this class:
their equality walls can be irrational. The construction below takes
maxima only of rational component functions, then adds them with real
weights.

Definitions 3.2--3.3: a \textbf{polyhedral convex combination} (PCC)
consists of a finite family \(x_i\in X\cap N_{\mathbb Q}\), real
coefficients \(a_i\in[0,1]\) with \(\sum_i a_i=1\), and centre
\(x=\sum_i a_i x_i\); on some rational structure of \(X\) there must be
cells \(\tau\preceq\sigma_i\) with \(x\in\tau\), \(x_i\in\sigma_i\). Raw
PC means that a finite real-valued function \(h\) satisfies

\[
h(x)\leq\sum_i a_i h(x_i)
\]

for every such PCC. We use neither the cited automatic continuity
theorem for raw PC functions nor a global ambient extension theorem. All
PA functions in this paper are continuous by definition.

Definition 3.1 tests PAPC by \(a\cdot Z\geq0\) against all nonnegative
local cycles of every positive dimension, on every ordinary open. We
prove the needed implication from PA PC in the real coefficient and real
weight scope, rather than assume the source's Lemma 3.4 extends
silently.

Definition 2.31 gives the divisor formula, writing \(a_\sigma\) for the
linear part of an affine piece:

\[
(a\cdot Z)(\tau)=
\sum_{\sigma\succ\tau} Z(\sigma)a_\sigma(n_{\sigma/\tau})
-a_\tau\left(\sum_{\sigma\succ\tau}Z(\sigma)n_{\sigma/\tau}\right).
\tag{D}
\]

The last term is the \textbf{tangential correction}. It will be nonzero
on the original rays in the decisive calculation. Remark 2.33 and Lemmas
2.34--2.35 state independence of normal lifts, refinement compatibility,
linearity, balancing, and commutativity. Their proof of balancing and
commutativity invokes Allermann--Rau Proposition 3.7. We do not import
that theorem: the complete two-dimensional algebra needed here is proved
below, including saturated plane indices.

Definition 4.4 defines the PA measure by the finite atomic zero-cycle
\(a^2\cdot[X]\) in our dimension two. This is the recursive intersection
normalization, without a division by \(2!\) or a probability
normalization. We use the explicit real-coefficient counterpart of this
definition. Positivity, finiteness, refinement independence, and
locality are proved below. No top-cell-interior Alexandrov formula is
substituted for a calculation on the lower faces.

Definition 3.14 requires \(a-\gamma=O(1)\) for the reference growth
class. Definition 4.15 requires positive reference slope on every
nonzero recession direction of every affine cell. Definitions 3.20, 3.23
and Proposition 3.22 describe regularizability through decreasing PAPC
sequences and uniform PAPC approximation. For our particular \(P\), the
sequence \(P+1/n\) supplies both definitions explicitly; constants have
zero divisors and the measures are identical. No general equivalence or
uniform operator-extension theorem is needed for this case. Any
regularized extension agreeing with the PA operator has this same value
at \(P\).

Definition 6.2 is the raw envelope

\[
P_\gamma(u)(x)=\sup\{h(x):h\in\mathrm{PC}(X),\ h\leq\gamma+u,
\ h\leq\gamma+A_h\text{ for some }A_h\in\mathbb R\}. \tag{E}
\]

Definition 6.9 first assumes pair regularity and then asks for the
contact integral for every bounded PA obstacle. Our negative witness has
an operator value in its own right, so it disproves the actual integral
identity without any assumption about other envelopes. Source Theorems
4.13 and 6.4, Proposition 4.16, analytic comparisons, and total-mass
theorems are not premises of the result.

The prior product paper, its embedded evidence and audit receipt, and
the six-cone parent proposal are fallible research data. We reconstruct
every lemma used from them. Their PASS labels are not proofs or
certificates for these new bytes. The book and primary example are used
for comparison only. No earlier local note is an unreproved premise of
this argument.

\subsection{Finite-dimensional
preliminaries}\label{finite-dimensional-preliminaries}

\textbf{Lemma 1 (polyhedral finite generation and closedness).} A
rational polyhedral cone has finitely many rational generators,
including both signs of a basis for its lineality. A nonempty pointed
rational polyhedron is \(\operatorname{conv}(A)+\operatorname{cone}(B)\)
with finite nonempty rational \(A\) and finite rational \(B\). The
projection of a finite nonstrict linear inequality system is again such
a system, retaining rationality for rational data. Every finitely
generated real cone is closed.

\textbf{Proof.} First consider a pointed cone written with finitely many
rational forms \(\ell_i(v)\geq0\), encoding equalities as pairs. Their
common kernel is its lineality and hence zero. The sum
\(\ell=\sum_i\ell_i\) is positive at every nonzero cone vector: all
terms are nonnegative there, and simultaneous vanishing would give a
vector in that kernel. The section \(\ell=1\) is bounded. Otherwise,
normalize an unbounded sequence and extract a convergent subsequence on
the unit sphere; its limit would be a nonzero cone vector with
\(\ell=0\). This section is a compact polyhedron.

A point in a positive-dimensional minimal face of a compact polyhedron
lies on a line in that face, with two terminal points in smaller faces.
Induct on dimension to express every point as a convex combination of
vertices. Vertices are obtained by full-rank subsets of the finite
active equations in the affine hull; there are finitely many, and they
are rational for rational equations. These vertices generate the cone.
For the zero cone the empty generating list suffices. For lineality,
quotient by its rational common kernel, apply the pointed result, lift
rational generators, and append both signs of a rational basis of the
kernel.

For a nonempty pointed polyhedron \(Q=\{x:Ax\geq b\}\), homogenize to
\(\{(y,t):Ay\geq bt,\ t\geq0\}\). Its lineality would have \(t=0\) and
give lineality of \(Q\), so the cone is pointed. Generators with
positive last coordinate normalize to rational points of \(Q\);
last-coordinate-zero generators are recession directions. The section
\(t=1\) is exactly the asserted convex hull plus cone, with at least one
point generator by nonemptiness.

For projection, eliminate one variable: each of its finitely many lower
bounds must be at most each upper bound, while inequalities independent
of that variable remain. A finite one-sided list of bounds is always
satisfiable. This is exact over the real numbers and preserves nonstrict
inequalities and rational coefficients. Repetition proves the projection
assertion. A cone generated by \(g_i\) is the projection of
\(x=\sum_i t_i g_i\), \(t_i\geq0\). The projection statement, valid also
for real coefficients, makes it a closed polyhedron. This proves the
last assertion. \(\square\)

\textbf{Lemma 2 (finite support alternative).} For finite vectors
\(d_i\) and real numbers \(b_i\), there is a linear form \(p\) with
\(p(d_i)\leq b_i\) for all \(i\) if and only if

\[
\sum_i t_i b_i\geq0\quad\text{whenever }t_i\geq0,
\quad \sum_i t_i d_i=0. \tag{F}
\]

\textbf{Proof.} Necessity follows by summation. For sufficiency consider
the closed cone \(K=\operatorname{cone}\{(d_i,b_i),(0,1)\}\), closed by
Lemma 1. Condition (F) excludes \(z=(0,-1)\) from \(K\). A nearest point
\(k\in K\) to \(z\) exists by restricting a minimizing sequence to a
large compact ball. Put \(r=k-z\). Differentiating the squared distance
along the segment from \(k\) to any \(w\in K\) gives
\(\langle r,w-k\rangle\geq0\). Taking \(w=0,2k\) gives
\(\langle r,k\rangle=0\). Hence \(\langle r,w\rangle\geq0\) on \(K\) and
\(\langle r,z\rangle=-\|r\|^2<0\). Write \(r=(-p_0,\alpha)\); the latter
inequality gives \(\alpha>0\). Evaluation on each generator gives
\(p_0(d_i)\leq\alpha b_i\). Take \(p=p_0/\alpha\). \(\square\)

\subsection{Actual PCC convexity on fixed rational
domains}\label{actual-pcc-convexity-on-fixed-rational-domains}

Let \(\Omega\) be a finite rational complex in a lattice vector space
\(W\), with carrier \(Y\). No balancing assumption is needed for the
next two lemmas. For a face \(\tau\preceq\sigma\) set

\[
T_{\sigma,\tau}=\{v:x+tv\in\sigma\text{ for all sufficiently small }t\geq0\},
\quad x\in\operatorname{relint}(\tau).
\]

This cone depends only on the inequalities active on \(\tau\), and not
on the chosen relative-interior point. It is rational, and has finite
rational generators \(d_{\sigma,\tau,r}\) by Lemma 1. Its lineality
contains the tangent space of \(\tau\).

\textbf{Lemma 3 (finite lifted local-support system).} A continuous real
PA function \(a\) on \(\Omega\) is PC on \(Y\) if and only if, for every
face \(\tau\), there is an ambient linear form \(p_\tau\in W^*\) such
that

\[
p_\tau(d_{\sigma,\tau,r})
\leq a_\sigma(d_{\sigma,\tau,r})
\quad\text{for every }\sigma\succeq\tau,r. \tag{L}
\]

Together with continuity of the affine pieces, these are finitely many
rational linear equations and inequalities in the affine coefficients of
\(a\) and the coefficients of \(p_\tau\). They characterize all actual
PCCs, even on other rational structures of \(Y\).

\textbf{Proof of necessity.} Choose
\(x\in\operatorname{relint}(\tau)\cap N_{\mathbb Q}\). A common small
rational \(t>0\) makes all \(x+td_i\) rational and places each in its
incident cell. For \(t_i\geq0\) with \(\sum_i t_i d_i=0\) and
\(T=\sum_i t_i>0\), these endpoints and real coefficients \(t_i/T\) form
a PCC with centre \(x\) on \(\Omega\). The PC inequality and affinity
imply \(\sum_i t_i a_{\sigma_i}(d_i)\geq0\). For \(T=0\) this is
automatic. Apply Lemma 2 to get \(p_\tau\).

\textbf{Proof of sufficiency.} Extend (L) to each tangent cone by
nonnegative linear combination. At every
\(x\in\operatorname{relint}(\tau)\), in each incident affine cell, this
gives

\[
a(x+v)\geq a(x)+p_\tau(v) \tag{G}
\]

for all sufficiently small feasible \(v\). A cell containing \(x\)
contains \(\tau\): its intersection with \(\tau\) is a face containing a
relative-interior point of \(\tau\). The finitely many cells not
containing \(x\) have positive distance from \(x\). Therefore (G) is a
simultaneous supporting inequality on the entire carrier germ, with one
\(p_\tau\).

On a segment \([x,y]\subset Y\), \(a\) restricts to a continuous PL
function with finitely many intervals. At a breakpoint \(z\), let \(p\)
be its germ support and \(v=y-x\). The left and right slopes \(b_-,b_+\)
satisfy \(b_-\leq p(v)\leq b_+\), using (G) in the negative and positive
directions. Thus slopes are nondecreasing. At the starting point the
initial slope is at least \(p_\tau(y-x)\), so integration along the
segment gives

\[
a(y)\geq a(x)+p_\tau(y-x). \tag{H}
\]

For any PCC on any rational structure, every centre-to-endpoint segment
lies in \(Y\), because its endpoint cell is convex and contains its
centre face. Apply (H) to all endpoints with the same support form at
the centre. Weighted summation cancels the linear terms. This gives
exactly the PCC inequality, at real centres and with real coefficients.
No intermediate centre outside \(Y\) is introduced.

Finally choose rational basepoints and rational tangent bases in all
cell affine hulls. An affine assignment is a finite coefficient vector;
equality on a common face is tested at its rational basepoint and
tangent basis, giving rational linear continuity equations. Every
derivative evaluation on a rational generator is rational-linear in
these variables. The ambient form evaluations are also rational-linear.
This proves the stated finite rational system. \(\square\)

\textbf{Lemma 4 (PA PC effectivity in the full real scope).} A
continuous real PA PC function on finite rational domains satisfies (D)
with nonnegative coefficients against every nonnegative local real
balanced cycle of every positive dimension, on every ordinary open set.

\textbf{Proof.} Refine the function and the tested cycle to a common
rational complex. At a tested \((k-1)\)-face \(\tau\), use Lemma 3 for
this refinement and a point in its relative interior. The two tangent
directions along \(\tau\) show that \(p_\tau=a_\tau\) on its tangent
space. Every inward primitive normal is in \(T_{\sigma,\tau}\): even if
a lift has a tangential component, a sufficiently small move from a
relative-interior point enters the incident cell. Consequently
\(a_\sigma(n)-p_\tau(n)\geq0\). Balancing and (D) give

\[
(a\cdot Z)(\tau)=\sum_{\sigma\succ\tau}Z(\sigma)
\bigl(a_\sigma-p_\tau\bigr)(n_{\sigma/\tau})\geq0.
\]

The same germs are tested after restriction to an ordinary open. This
treats every dimension \(1\leq k\leq\dim Y\), every local cycle
structure, and arbitrary real nonnegative weights. \(\square\)

This is sufficient for all PAPC assertions below. No equivalence for
non-PA raw functions or global approximation theorem has been inferred.

\subsection{Exact rational decomposition of every real PA
reference}\label{exact-rational-decomposition-of-every-real-pa-reference}

\textbf{Lemma 5 (exact finite convex decomposition with strict recession
growth).} If \(\gamma\in\mathrm{PA}_{\mathbb R}(Y)\cap\mathrm{PC}(Y)\)
on a fixed finite rational complex is positive on every nonzero
recession direction of every affine cell, then

\[
\gamma=\sum_{j=1}^s a_j\gamma_j,\qquad a_j>0,\quad\sum_j a_j=1, \tag{R}
\]

where every \(\gamma_j\in\mathrm{PA}_{\mathbb Q}(Y)\cap\mathrm{PC}(Y)\)
has the same strict growth property on that same complex. Equality is
exact on the entire carrier, including every unbounded cell.

\textbf{Proof.} Lemma 3 supplies a feasible lifted coefficient vector
\(y\), consisting of the affine pieces of \(\gamma\) and their support
forms. It satisfies rational homogeneous equations \(Ey=0\) and rational
nonstrict inequalities \(Ay\geq0\).

Each reference cell is pointed: nonzero lineality would give opposite
nonzero recession directions with opposite slopes, both required to be
positive. By Lemma 1 its recession cone has finitely many nonzero
rational generators. Strict growth is precisely a finite list \(By>0\)
of rational-linear slope inequalities on these generators. Positive
values on them imply positivity on every nonzero nonnegative
combination; bounded cells impose no condition.

Adjoin to \(E\) the rows of \(A\) that vanish at \(y\), and let \(L\) be
the resulting kernel. This is a rational linear subspace and contains
\(y\). In \(L\), all remaining rows of \(A\) and all rows of \(B\) are
strictly positive near \(y\). Choose a rational basis of \(L\), and a
sufficiently small coordinate box with rational endpoints strictly
enclosing the coordinates of \(y\). Every point of the box satisfies
those strict inequalities; all active rows remain zero. Each box vertex
\(y_j\) is rational, continuous, locally support-feasible, and strictly
growing.

In one coordinate \(l<t<h\), the weights \((h-t)/(h-l)\),
\((t-l)/(h-l)\) are positive and reproduce \(t\). Products of these
weights express \(y\) exactly as a convex combination of the finitely
many box vertices. If \(L\) is zero-dimensional, its only vector is
rational and a singleton representation suffices when feasible.
Projecting the exact identity onto the affine pieces gives (R). Lemma 3
proves each projected function PC. The finite recession tests prove its
growth. \(\square\)

This reconstructs the fixed-domain decomposition suggested by the prior
product paper. It is not an approximation of a reference by
bounded-difference rational references: the individual \(\gamma_j\) can
have different recession slopes. Their original real-weighted sum is
exactly \(\gamma\), and that is what will restore the exterior reference
below. No ambient convex extension of \(\gamma\) is asserted.

\subsection{Recursive real divisor algebra in dimension
two}\label{recursive-real-divisor-algebra-in-dimension-two}

Formula (D) is bilinear in the affine coefficients and cycle weights. A
change of normal lift lies in \(N_\tau\); the incident slopes restrict
to \(a_\tau\) there, so its contribution cancels between the two terms
of (D). An ambient affine function has zero divisor on every balanced
cycle, because its common linear part factors out of the sum. Thus
constants and ambient affine functions may be subtracted at each
recursive step. The formula uses only incident germs, proving locality.

For refinement, fix a new \((k-1)\)-face and its minimal old face. If
the old face has dimension \(k-1\), the saturated tangents, inherited
incident weights, and quotient normal classes reproduce the old
coefficient. If it has dimension \(k\), a generic point of the new face
lies inside an old \(k\)-cell; there are two local subdivision sides
with equal inherited weights and slopes and opposite normals, so the
coefficient is zero. If the old face has dimension greater than \(k\),
no supported old \(k\)-cell can contain the new face, so its coefficient
is zero. These cases exhaust the possibilities. This proves refinement
compatibility both for cycles and their divisors, locally on ordinary
opens as well.

For balancing and commutativity in dimension two, work at a vertex germ.
Refine its rational local fan into pointed two-cones if it originally
has lineality. Every such cone \(C_{rs}\) has primitive boundary rays
\(u_r,u_s\) and saturated plane index

\[
\chi_{rs}=[N_{C_{rs}}:\mathbb Z u_r+\mathbb Z u_s].
\]

If its weight is \(\omega_{rs}\), an inward primitive normal along \(r\)
equals \(u_s/\chi_{rs}\) modulo \(\mathbb R u_r\). To see this, extend
\(u_r\) to a basis of its saturated plane lattice. The second coordinate
of \(u_s\) is, up to the inward orientation, \(\chi_{rs}\). The second
basis vector is an integral primitive normal, differing from
\(u_s/\chi_{rs}\) by a real tangent vector. The cancellation just proved
allows this latter expression in (D).

Balancing of the original two-cycle at \(r\) therefore says

\[
\sum_{C_{rs}\succ r}\frac{\omega_{rs}}{\chi_{rs}}u_s=t_r u_r
\]

for a real \(t_r\). Let (a(u)) denote the value of the linear germ on a
ray generator, after subtracting the vertex value. Its first divisor
coefficient is

\[
D_a(r)=\sum_{C_{rs}\succ r}\frac{\omega_{rs}}{\chi_{rs}}a(u_s)-t_r a(u_r). \tag{A}
\]

Summing \(D_a(r)u_r\), the first term contributes, per unordered cone,

\[
\frac{\omega_{rs}}{\chi_{rs}}
\bigl(a(u_s)u_r+a(u_r)u_s\bigr).
\]

Substitution of the displayed balancing relation into the second term
contributes exactly the same expression. Thus \(\sum_r D_a(r)u_r=0\):
the first divisor is a balanced one-cycle. This local argument proves
balancing at every vertex of every finite two-cycle. A divisor of a
one-cycle is a zero-cycle and has no further balancing condition.

For a second function \(b\), the second divisor at the vertex is

\[
\sum_{\{r,s\}}\frac{\omega_{rs}}{\chi_{rs}}
\bigl(a(u_s)b(u_r)+a(u_r)b(u_s)\bigr)
-\sum_r t_r a(u_r)b(u_r). \tag{B}
\]

It is symmetric in (a,b), proving commutativity. This proof works
directly for real coefficients and real weights; no rational theorem is
being invoked by density. It applies at every vertex and is compatible
with the refinements already checked.

The same calculation proves the corresponding algebra for the auxiliary
local-cycle tests in arbitrary dimension. At a codimension-two face of a
k-cycle, take its local fan and quotient by that face's saturated
tangent lattice. Every supported local k-cone contains this tangent
space as lineality and becomes a two-cone; every incident (k-1)-face
becomes a ray. Subtract a common ambient linear extension of each
function's slope along the face, so the residual functions factor
through the quotient. The quotient tangent lattices are saturated: if an
integer multiple of a class lies in N\_sigma/N\_eta, then an integer
multiple of its lift lies in N\_sigma, and saturation of N\_sigma
implies that its lift already lies there. Primitive quotient normals
therefore coincide with the images of the old normals, with exactly the
plane indices used above. Balancing of the old cycle becomes balancing
of this two-cycle. Formula (D) passes to this quotient because the
removed tangential linear part has zero divisor. The two-dimensional
cancellation consequently proves balancing of the first divisor modulo
N\_eta, which is precisely balancing along the original tested face, and
the symmetric expression proves commutativity there. Testing all
codimension-two faces proves the assertion; k=1 has a zero-cycle divisor
and no balancing condition. This is local quotient algebra, not a
cycle-pushforward theorem.

For every PA PC \(a\) on \(X\), Lemma 4 first makes \(a\cdot[X]\)
nonnegative, and then makes \(a\cdot(a\cdot[X])\) nonnegative, since the
first divisor is an actual local balanced one-cycle. A common finite
rational structure has finitely many vertices, so this is a finite
positive atomic measure. No boundary atom outside \(X\) is introduced by
the PA definition. This constructs precisely the measure needed in (T),
including all lower-face and exterior coefficients.

\subsection{The carrier, primitive normals, and the complete conical
mass}\label{the-carrier-primitive-normals-and-the-complete-conical-mass}

Any three distinct \(v_i\) are a \(\mathbb Z\)-basis of \(N\):
determinants are \(1\) or \(-1\). Thus every pair spans its saturated
lattice \(\mathbb Z v_i+\mathbb Z v_j\); each original ray has saturated
lattice \(\mathbb Z v_i\). The only real relation among the four
generators is \(v_1+v_2+v_3+v_4=0\).

The cones meet face to face. If two cone representations agree, their
four coefficient differences form a multiple of \((1,1,1,1)\). An index
missing from their union forces that multiple to be zero. If their union
has all four indices, the two pairs are disjoint and their differences
have opposing signs, again forcing the multiple to be zero. Their
intersection is consequently their common ray when a generator is
shared, or the origin otherwise. This proves that \(X\) is a finite
rational pure two-dimensional carrier.

Along \(\mathbb R_{\geq0}v_i\), use \(v_j\) as the inward primitive
normal lift from each incident \(\sigma_{ij}\). Saturation makes its
quotient class primitive; adding a large multiple of \(v_i\) produces a
lift geometrically entering the cone if desired. Their sum is \(-v_i\),
tangent to the ray. Hence the six unit weights balance. Refinement
preserves this balancing by the proved cancellation rule.

Both \(q\) and \(\rho\) are restrictions of continuous convex ambient
rational PL functions, and satisfy all PCC inequalities by ordinary
Jensen. They are therefore PA PC and PAPC by Lemma 4. The norm \(\rho\)
is proper on this closed carrier, positive away from zero, with
\(0\leq q\leq\rho\). On a refined cell where \(\rho\) is affine, its
slope on any nonzero recession direction \(v\) is \(\rho(v)>0\). Indeed
the sign of each coordinate along an unbounded cell is constant and
agrees with the sign of its recession coordinate; summing gives that
slope.

In parameters \(x=t v_i+s v_j\), \(t,s\geq0\), the full table of \(q\)
is

\begin{longtable}[]{@{}ll@{}}
\toprule\noalign{}
Cone & \(q(t v_i+s v_j)\) \\
\midrule\noalign{}
\endhead
\bottomrule\noalign{}
\endlastfoot
\(12,34\) & \(\lvert t-s\rvert\) \\
\(13,14,23,24\) & \(t+s\) \\
\end{longtable}

Thus the only interior bends of \(q\) are on \(h_{12}=v_1+v_2=(1,1,0)\)
and \(h_{34}=v_3+v_4=(-1,-1,0)\). Both are primitive in their actual
saturated planes. At either diagonal choose \(v_i\) and \(v_j\) as the
respective inward primitive normals of the two refined halves. Their sum
is \(h_{ij}\), and the tangent derivative of \(q\) on \(h_{ij}\) is
zero. The two inward derivatives are \(1,1\), so (D) gives divisor
weight two.

Along an original ray \(v_i\), the linear germ in the incident cone
paired with its exceptional partner has inward derivative \(-1\), and
the other two incident germs have derivatives \(1,1\). These are the
derivatives for the chosen normal lifts \(v_j\), also after diagonal
refinement: near \(v_i\), a move in the \(v_j\) direction stays on the
relevant half for sufficiently small displacement. All four tangent
slopes are \(q(v_i)=1\). The sum of the normal lifts is \(-v_i\), so the
correction in (D) is \(-1\). Consequently the original-ray weight is
\(1-(-1)=2\), not the uncorrected sum one.

There are no other first-divisor faces: on each remaining
relative-interior affine sector the function has one linear piece, and
any artificial refinement face cancels. The complete first divisor is

\[
q\cdot[X]=2\sum_{i=1}^4[\mathbb R_{\geq0}v_i]
+2[\mathbb R_{\geq0}h_{12}]+2[\mathbb R_{\geq0}h_{34}]. \tag{Q1}
\]

It is balanced directly because \(\sum_i v_i=0\) and
\(h_{12}+h_{34}=0\). On this one-cycle, \(q\) has slope one on the four
original rays and slope zero on the two diagonal rays. The only vertex
is the origin; on every open ray there is no new bend. Therefore the
complete second divisor is

\[
q^2\cdot[X]=\left(\sum_{i=1}^4 2\cdot1+2\cdot0+2\cdot0\right)\delta_0
=8\delta_0. \tag{Q2}
\]

This computes the actual primitive-normal recursive intersection. It
uses no difference-map pushforward, proper or nonproper, and includes
the original-ray tangential corrections and the lower-dimensional
diagonal contributions. The number eight is not an Alexandrov mass
divided by a factorial.

To identify the named carrier without a smoothness inference, use the
lattice \(\mathbb Z^4/\mathbb Z(1,1,1,1)\simeq N\), sending the four
coordinate classes to \(v_i\). The uniform matroid \(U_{3,4}\) has the
sole circuit \(\{1,2,3,4\}\). Its circuit-definition Bergman carrier
consists of classes whose minimum among the four coordinates is achieved
at least twice. Subtracting that minimum yields a representative with
nonnegative coordinates and at most two positive coordinates. Its image
is exactly one of the six cones above. The flag-of-flats fan subdivides
each \(\sigma_{ij}\) by \(v_i+v_j\): the proper nonempty flats are
singletons and pairs, and a singleton contained in a pair gives the two
half-cones. This also checks that the unit-weight carrier is the one
displayed in BMPS Example 6.24. No undefined higher-dimensional BMPS
polyhedral-smoothness condition is asserted.

\subsection{Gluing to every real reference, with exact exterior
equality}\label{gluing-to-every-real-reference-with-exact-exterior-equality}

Fix an arbitrary \(\gamma\) from the theorem and a rational \(R>0\), for
example \(R=1\). Refine its finite rational affine structure by the
domains of \(\rho\), if necessary. Strict recession positivity is
inherited by every refined cell. Apply Lemma 5 to get the exact
decomposition (R), with all \(\gamma_j\) rational PA PC and strictly
growing.

For each reference cell, Lemma 1 expresses it as a convex hull of
finitely many rational points plus a finitely generated rational
recession cone. At every nonzero recession generator, both the \(\rho\)
slope and each \(\gamma_j\) slope are positive. Choose a common rational
\(\eta>0\) smaller than every ratio of these latter slopes to the
former. There are finitely many ratios. On each cell
\(\gamma_j-\eta\rho\) is affine with nonnegative recession slopes, so
its minimum is bounded below by the finitely many point-generator
values. Choose one positive rational \(B\) larger than all negative
minima. Thus

\[
\gamma_j\geq\eta\rho-B\quad\text{on }X\quad(1\leq j\leq s). \tag{G1}
\]

Bounded cells are handled by their point generators and add no recession
conditions. This proves a quantitative bound, not just a claim of
coercivity.

Choose a common rational \(0<\lambda<\eta\). On the compact set
\(\rho\leq2R\), each continuous \(\gamma_j-\lambda q\) is bounded above.
Choose positive rational \(M_j\) strictly larger than that upper bound.
Define

\[
P_j=\max\{\lambda q,\gamma_j-M_j\},\quad
P_0=\sum_j a_j P_j,\quad C=\sum_j a_j M_j,\quad P=P_0+C. \tag{G2}
\]

Finite maxima preserve actual PCC inequalities: choose a maximizing
function at the centre, apply its PC inequality, and then bound its
endpoint values by the endpoint maxima. Nonnegative sums preserve them
as well. Therefore \(P_j,P_0,P\) are PC. Each maximum in (G2) is between
rational affine pieces, so its comparison walls and spatial subdivision
are rational. A finite rational common refinement makes the
real-weighted sum \(P_0\), and then \(P\), real PA. Adding real
coefficients creates no new walls. No direct maximum of \(\gamma\) with
a rational function is declared spatially rational.

Every \(P_j\geq\gamma_j-M_j\), so \(P_0\geq\gamma-C\) and
\(P\geq\gamma\). Also

\[
\lambda q-\gamma_j\leq(\lambda-\eta)\rho+B,
\]

which is bounded above. Therefore \(P_j-\gamma_j\), \(P_0-\gamma\), and
\(P-\gamma\) are bounded. Lemma 4 gives full local-cycle PAPC
effectivity for \(P\).

By construction \(P_j=\lambda q\) on \(\rho\leq2R\), hence
\(P=C+\lambda q\) there. From (G1) and \(q\leq\rho\),

\[
\gamma_j-M_j-\lambda q\geq(\eta-\lambda)\rho-B-M_j.
\]

Choose a common rational \(S>2R\) with
\((\eta-\lambda)S\geq B+\max_j M_j\). Then on \(\rho\geq S\) every
\(P_j=\gamma_j-M_j\), and the exact original weights restore

\[
P=P_0+C=\gamma. \tag{G3}
\]

This is exact exterior equality with the prescribed reference, not
merely equality of asymptotic slopes or bounded-difference
approximation. The construction is finite and existential for an
arbitrary real coefficient input; it does not posit an algorithm
deciding exact equalities of arbitrary computably unspecified real
numbers.

The sequence \(P+1/n\) consists of real PA PC functions with the same
bounded reference difference, decreases uniformly to \(P\), and has the
same recursive divisor measure at every \(n\). It proves membership of
the particular \(P\) in the real regularizable domain. We have not
assumed or concluded whole-pair regularity.

\subsection{Raw envelope proof at every real
point}\label{raw-envelope-proof-at-every-real-point}

Set

\[
b_R=(1-\rho/R)_+,\qquad f=P+\epsilon b_R,\qquad u=f-\gamma, \tag{O}
\]

where \(\epsilon\in\{1,2\}\) will be fixed below. The roof is rational
PA, supported on the compact set \(\rho\leq R\). By (G3), \(u=0\) on
\(\rho\geq S\), and \(u\geq0\) everywhere because \(P\geq\gamma\). It is
a bounded continuous real PA obstacle on finite rational spatial
domains.

The rational points

\[
y_+=\frac R2 h_{12},\qquad y_-=\frac R2 h_{34}=-y_+
\]

have \(\rho=R\), \(q=0\), and \(f=P=C\). With coefficients \(1/2,1/2\)
and origin centre, they form an actual PCC: the face \(\{0\}\) is
incident to the endpoint cones \(\sigma_{12},\sigma_{34}\). Thus every
raw PC \(h\leq f\), including every competitor in (E), satisfies

\[
h(0)\leq C=P(0). \tag{E0}
\]

For a rational \(x\in X\) with \(0<\rho(x)<R\), the rational point
\(y=Rx/\rho(x)\) is in the same closed original cone and has
\(\rho(y)=R\). The endpoints \(0,y\), centre \(x\), and coefficients
\(1-\rho(x)/R,\rho(x)/R\) form a same-cell PCC on the original six-cone
structure. Endpoints need not be in relative interiors under Definition
3.2. The structure of the obstacle need not be this PCC structure. Since
\(f(y)=C+\lambda q(y)\), (E0) and homogeneity give

\[
h(x)\leq C+\frac{\rho(x)}R\lambda q(y)
=C+\lambda q(x)=P(x). \tag{E1}
\]

For an arbitrary real \(x\) with \(\rho(x)<R\), take its minimal face in
a finite rational common affine-domain structure for \(q,\rho\). If this
face is zero-dimensional, \(x\) is rational. Otherwise, in its rational
affine hull choose a small enclosing simplex with \(x\) strictly inside
and vertices in the relative interior of the face and in \(\rho<R\).
Perturb its vertices to rational points; for sufficiently small
perturbation the vertices stay in that face and ball and their
barycentric coordinates of \(x\) stay positive. This follows directly
from the continuous inverse of the affine-coordinate matrix of a
nondegenerate simplex. All vertices are now rational and have the bounds
(E0)--(E1). With their real barycentric coefficients they form a
same-cell PCC. Since \(C+\lambda q\) is affine there, its weighted
vertex values equal its value at \(x\), so \(h(x)\leq P(x)\). No
continuity, semicontinuity, approximation, or affine extendibility of
raw \(h\) has been used.

For \(\rho\geq R\), the roof is zero and \(h\leq f=P\) directly. Thus
every competitor in (E) lies below \(P\) on all \(X\). Conversely \(P\)
is PC, lies below \(f\), and has bounded reference difference, so is
itself a competitor. We conclude

\[
P_\gamma(u)=P\quad\text{everywhere on }X. \tag{EP}
\]

In particular, there is no remote growth loophole: the lower competitor
exists globally with the prescribed exact exterior equality, and the
upper bounds apply to all raw competitors, even without their growth
restriction.

For nonseparability choose the four rational points
\(0,Rv_1,Rv_3,R(v_1+v_3)\) in \(\sigma_{13}\). Their \(\rho\) values are
\(0,R,R,2R\); define

\[
\Delta a=a(0)+a(R(v_1+v_3))-a(Rv_1)-a(Rv_3).
\]

A sum of a function of the first cone parameter and a function of the
second has \(\Delta a=0\). Here \(\Delta b_R=1\). Let
\(A=\Delta(P-\gamma)\) and choose \(\epsilon=1\) unless \(A=-1\), in
which case choose \(\epsilon=2\). Then \(\Delta u=A+\epsilon\ne0\). This
choice changes neither the endpoint equalities nor the envelope
argument, and gives the promised genuinely nonseparable obstacle
restriction.

\subsection{Complete contact defect and operator
domain}\label{complete-contact-defect-and-operator-domain}

Take any finite rational common refinement \(\Omega\) for \(P,[X]\) and
the roof. By the recursive algebra and Lemma 4,

\[
\mu=\mathrm{MA}_{\mathrm{poly},\mathbb R}(P)
=\sum_{z\in\Omega^{(0)}}m_z\delta_z,
\qquad m_z=(P\cdot(P\cdot[X]))(z)\geq0
\]

is a well-defined finite positive measure, independent of refinement.
Every \(m_z\) uses all incident first-divisor edges through (D); lower
faces are not omitted. The same coefficients are obtained by the
decreasing regularizing sequence \(P+1/n\).

On the ordinary open \(\rho<2R\), the exact equality \(P=C+\lambda q\),
constant invariance and recursive locality give

\[
\mathbf1_{\{\rho<2R\}}\mu=8\lambda^2\delta_0. \tag{LMA}
\]

In particular this holds on the roof's nonzero set \(\rho<R\). The full
integrand is exactly \(P-f=-\epsilon b_R\). At the origin it is
\(-\epsilon\), and it is zero at every other atom: in \(0<\rho<R\) there
is no mass by (LMA); on \(\rho\geq R\), including the roof boundary, the
interpolation region, all lower faces, every exterior vertex, and the
exact-growth region, the integrand itself is zero. Thus summing the
entire finite measure, not just its central atom and not just top-cell
interiors, yields (T).

This proves the full operator value and contact defect for the
constructed envelope, without invoking whole-pair regularity, total
mass, or an extension theorem for non-PA functions. The outside
coefficients may depend on \(\gamma\); their contribution to this exact
integral is rigorously zero. No compactification measure is being
inserted, and no analytic realization or literal coefficient conclusion
follows.

\subsection{Explicit reference: full first divisor and all nine
atoms}\label{explicit-reference-full-first-divisor-and-all-nine-atoms}

For \(\gamma=\rho\), choose \(R=1\), \(\lambda=1/4\), \(M=4\), and
\(\epsilon=1\). It is useful first to write the unshifted candidate

\[
P_0=\max\{q/4,\rho-4\},\quad b=(1-\rho)_+.
\]

Then the theorem's actual envelope is \(P=P_0+4=\max\{4+q/4,\rho\}\) and
\(u=P+b-\rho\). Translation does not change its measure or defect. On
\(\rho\leq2\), \(\rho-4\leq-2\leq q/4\), so \(P_0=q/4\). Since
\(q\leq\rho\), on \(\rho\geq16/3\) we have \(\rho-4\geq q/4\). Thus
\(P=\rho\) there and \(u\) is compactly supported. On \(\sigma_{13}\),
\(\Delta u=1\), so this explicit obstacle is nonseparable too.

The following direct table also computes all exterior mass. Put

\[
A_i=(16/3)v_i\ (i=1,2,3),\quad A_4=(16/11)v_4,
\]

\[
D_+=2h_{12},\quad D_-=2h_{34},\quad
D_{14}=(8/3)(v_1+v_4),\quad D_{24}=(8/3)(v_2+v_4).
\]

Using (D), the \textbf{complete nonzero first divisor}
\(P\cdot[X]=P_0\cdot[X]\) has the following 22 edges. A closed segment
denotes its relative-interior weight, and a continuing ray starts at its
indicated endpoint in the indicated direction.

\begin{longtable}[]{@{}lr@{}}
\toprule\noalign{}
Edges & Weight \\
\midrule\noalign{}
\endhead
\bottomrule\noalign{}
\endlastfoot
\([0,A_i]\), \(i=1,2,3\) & \(1/2\) \\
Rays from \(A_i\) in direction \(v_i\), \(i=1,2,3\) & \(4\) \\
\([0,A_4]\) & \(1/2\) \\
\([0,D_+]\), \([0,D_-]\) & \(1/2\) \\
Ray from \(D_-\) in direction \(h_{34}\) & \(2\) \\
Rays from \(D_{14},D_{24}\) in directions \(v_1+v_4,v_2+v_4\) & \(2\) \\
\([A_1,D_+]\), \([D_+,A_2]\) & \(1/4\) \\
\([A_1,A_3]\), \([A_2,A_3]\) & \(3/4\) \\
\([A_4,D_{14}]\), \([A_4,D_{24}]\) & \(1/4\) \\
\([D_{14},A_1]\), \([D_{24},A_2]\) & \(3/4\) \\
\([A_4,D_-]\), \([D_-,A_3]\) & \(1/4\) \\
\end{longtable}

For completeness, this table follows as follows. On cones \(12,13,23\),
\(\rho=t+s\). On \(14,24,34\), \(\rho=t+s\) for \(t\geq s\) and
\(\rho=3s-t\) for \(t\leq s\). Combine these with the preceding \(q\)
table. The maximum's comparison edges are:

\begin{longtable}[]{@{}llr@{}}
\toprule\noalign{}
Cone and sector & Comparison equation & Weight \\
\midrule\noalign{}
\endhead
\bottomrule\noalign{}
\endlastfoot
\(12, t\geq s\) & \(3t+5s=16\) & \(1/4\) \\
\(12, t\leq s\) & \(5t+3s=16\) & \(1/4\) \\
\(13,23\) & \(3t+3s=16\) & \(3/4\) \\
\(14,24, t\geq s\) & \(3t+3s=16\) & \(3/4\) \\
\(14,24, t\leq s\) & \(-5t+11s=16\) & \(1/4\) \\
\(34, t\geq s\) & \(3t+5s=16\) & \(1/4\) \\
\(34, t\leq s\) & \(-3t+11s=16\) & \(1/4\) \\
\end{longtable}

On each saturated parameter lattice \(\mathbb Z^2\), the difference of
the two competing gradients is one fourth of the displayed integral
normal covector. Dividing that covector by the gcd of its two entries
makes it primitive. Bezout supplies an inward normal with primitive
covector value one; the bend weight is its gradient difference, namely
that gcd divided by four. This proves every comparison-edge weight and
position. Original rays switch at the \(A_i\). Near them the central
\(q/4\) divisor has weight \(1/2\). Outside, the \(\rho\) divisor has
weights \(4,4,4,0\) on \(v_1,v_2,v_3,v_4\): all three inward derivatives
at \(v_i\), \(i\leq3\), are one and their tangential correction is
\(-1\); at \(v_4\) they are \(-1,-1,-1\) with correction \(-3\). The
only \(\rho\) interior bends are the three \(v_i+v_4\) diagonals, each
with weight two, since the gradient difference is twice the primitive
covector \(t-s\). The \(q\) diagonals inside its active region have
weights \(1/2\). Their switching positions are precisely
\(D_+,D_-,D_{14},D_{24}\). All other sectors have a single affine piece,
so there are no further first-divisor edges.

Intersecting once more at each of the nine endpoint vertices gives:

\begin{longtable}[]{@{}
  >{\raggedright\arraybackslash}p{(\columnwidth - 4\tabcolsep) * \real{0.3000}}
  >{\raggedright\arraybackslash}p{(\columnwidth - 4\tabcolsep) * \real{0.3000}}
  >{\raggedleft\arraybackslash}p{(\columnwidth - 4\tabcolsep) * \real{0.4000}}@{}}
\toprule\noalign{}
\begin{minipage}[b]{\linewidth}\raggedright
Vertex
\end{minipage} & \begin{minipage}[b]{\linewidth}\raggedright
Contributions
\(\text{edge weight}\times\text{primitive outgoing slope}\)
\end{minipage} & \begin{minipage}[b]{\linewidth}\raggedleft
Mass
\end{minipage} \\
\midrule\noalign{}
\endhead
\bottomrule\noalign{}
\endlastfoot
\(0\) & \(4(1/2)(1/4)\); the diagonal slopes are zero & \(1/2\) \\
\(A_1,A_2,A_3\), each & \(4-1/8-1/2\), all other terms zero &
\(27/8\) \\
\(A_4\) & \(-1/8+1+1-1/2\) & \(11/8\) \\
\(D_+\) & \(1/2+1/2\) & \(1\) \\
\(D_-\) & \(4+1/2+1/2\) & \(5\) \\
\(D_{14},D_{24}\), each & \(4-1+0\) & \(3\) \\
\end{longtable}

Here negative summands are legitimate individual directional slopes; the
complete coefficient is nonnegative. To check the non-axis slopes, at
\(A_1\) the direction toward \(D_+\) is \((-5,3,0)\), on which \(q/4\)
has outgoing slope \(-2\); its weight \(1/4\) gives \(-1/2\). At \(A_2\)
exchange (x,y). At \(A_3\) the direction toward \(D_-\) is
\((-3,-3,-8)\) with slope \(-2\). The other comparison slopes at these
vertices are zero. At \(A_4\), the directions toward \(D_{14},D_{24}\)
are \((6,-5,-5),(-5,6,-5)\), each with slope four; toward \(D_-\) it is
\((-3,-3,8)\) with slope \(-2\). At \(D_+\), the directions toward
\(A_1,A_2\) have slopes two. At \(D_-\), its outward diagonal slope is
two, and its directions toward \(A_4,A_3\) have slopes two. At
\(D_{14}\), the continuing diagonal has slope two, the direction toward
\(A_4\) has slope \(-4\), and the direction toward \(A_1\) has slope
zero. The \(D_{24}\) calculation exchanges (x,y). All these vectors are
primitive in the actual one-dimensional tangent lattices. This derives
the full recursive table rather than assuming it from numerical output.

The total mass is

\[
\frac12+3\frac{27}{8}+\frac{11}{8}+1+5+3+3=24.
\]

Every vertex except zero has \(\rho\geq4\), so its integrand is zero. At
zero the integrand is \(-1\). The complete integral is therefore
\(-1/2\). The total mass is an additional checked value, not the
argument for contact failure. Small exact-arithmetic diagnostics
independently reproduce all six saturated planes, all original-ray
corrections, all 22 first edges, all nine atoms and their weights, and
the same integral.

\subsection{Known comparisons and residual
questions}\label{known-comparisons-and-residual-questions}

BMPS v1 Section 6.5, Example 6.24 and Claim 6.25 already use
\textbf{this same carrier}, refined by all six \(v_i+v_j\) rays. Their
reference takes value one on all fine primitive generators except
\(h_{12},h_{34}\), on which it is zero. They construct an obstacle whose
envelope is that reference, with a noncontact atom of mass four. Its
zero line prevents positive growth at infinity. Those statements are
source comparisons, not imported inputs for our different function \(q\)
or its mass eight. We inspected the complete example and cone-by-cone
claim proof; the present recursive and envelope proofs do not rely on
either claim.

The controlling book Chapter 72 proposition
\texttt{prop:frontier-}\allowbreak{}\texttt{polyhedral-}\allowbreak{}\texttt{bergman-}\allowbreak{}\texttt{positive-}\allowbreak{}\texttt{infinity-}\allowbreak{}\texttt{obstruction},
including its full displayed proof, is snapshotted and was read. It
already repairs growth for a special reference by a maximum with two
shifted affine functions. Its conclusion is the dichotomy: the chosen
envelope is nonregularizable, \textbf{or}, if its operator is defined,
its contact integral is at most \(-4\varepsilon\). We do not count the
carrier, the existence of some positive-growth package failure, or that
special growth repair as new.

The result here strengthens that precise comparison in two ways: it
constructs an actual finite rational-domain PAPC envelope and its exact
full defect, without a regularity alternative, and it quantifies over
\textbf{every} positive-growth continuous real PA PC reference on this
carrier. The optional constant translation gives exact exterior equality
\(P=\gamma\) and a compactly supported nonnegative obstacle. The
specific \(\gamma=\rho\) illustration is evidence of the construction;
the all-reference theorem rests on the complete fixed-domain
decomposition and gluing proofs above.

The prior product paper proves an exact-scope criterion on
direct-lattice products of rank-one ray stars times affine space, using
the finite local-support and exact rational-reference decomposition. Its
review and read attestation are data. That carrier class is not assumed
here; the present six-cone carrier has its own proof. The reused
elementary lemmas have been fully reconstructed. No positive product
theorem, imported note n100, or prior PASS is a supporting premise.

The catalogue also records complete affine and finite graph results. We
do not repackage those as new progress. The original residual
classification concerns arbitrary higher-dimensional balanced rational
carriers and all admissible obstacles. We settle only the stated
universal-reference obstruction on the six-cone carrier. We do not
determine all positive carriers, prove arbitrary non-PA approximation,
or establish whole-pair regularity even on this negative carrier.
Arbitrary graph factors, other Bergman carriers, general intrinsic
contact criteria, and any precisely formulated higher-dimensional
smoothness/strict-positivity theorem remain outside this result.

BMPS Definition 6.19 defines polyhedral smoothness only in dimension
one. We make no higher-dimensional BMPS polyhedral-smoothness claim. The
circuit/flat calculation above establishes the named carrier
identification, not an undefined extra smoothness package. We also make
no PAPC+ claim under Definition 3.8; testing every strict PCC, and
resolving the literal allowance of zero coefficients in that wording,
would be a separate issue. An optional first-divisor strictness
calculation alone would not justify such a claim.

There is no literal integral-slope inference: the rational scale
\(\lambda\), rational component slopes, real convex weights and real
constant \(C\) need not preserve integral slopes or rational constants.
There is no analytic realization, CLD/delta-form/BFJ comparison,
polarized lift, or P05.5 conclusion. P06.2 has not been investigated in
this assignment. Status is \textbf{partial} and candidate is
\textbf{false}, because the original P05.2 classification remains open.

\subsection{Reproduction, evidence, and
bibliography}\label{reproduction-evidence-and-bibliography}

The authoritative complete argument is \texttt{P05.2-argument.md}. Its
standalone typeset versions are \texttt{P05.2-partial.tex} and
\texttt{build/P05.2-partial.pdf}. Run \texttt{bash\ rerun.sh} in this
fresh workspace. It executes \texttt{python3\ diagnostics.py}, creates
standalone TeX using installed Pandoc, compiles twice with installed
pdfLaTeX, and runs artifact checks. Exact diagnostic output is retained
in \texttt{diagnostic-output.json}; commands and logs are in
\texttt{build/}. \texttt{environment.json},
\texttt{source-manifest.json}, \texttt{artifact-hashes.json},
\texttt{author-self-audit.md}, \texttt{research-notes.md}, and the
five-key \texttt{status.json} document the environment, immutable source
binding, build outputs, author audit, scope and remaining gaps.

The finite diagnostics are exhaustive only for the explicitly enumerated
cells and edges of the displayed rational example. They are not
exhaustive over references, are not proof verification, and do not
certify the universal assertions. Two successful LaTeX passes establish
typesetting only. Hashes bind bytes, not mathematical correctness.

No installation, additional model call, collaboration subagent, heavy
computation, network retrieval, external message, authorization/proxy/CA
change, SQL access or recovery, or old-artifact write was used. No
operator-provided supervised heavy-compute facility was supplied; none
was needed for this finite calculation. All new artifacts are in the
assigned fresh workspace.

\begin{enumerate}
\def\labelenumi{\arabic{enumi}.}
\tightlist
\item
  A. M. Botero, E. Mazzon and L. Pille-Schneider, \emph{Monge--Ampère
  measures on balanced polyhedral spaces}, arXiv:2603.08664v1 (2026).
  Exact numbered definitions used above; Example 6.24/Claim 6.25 and
  Remark 6.26 are known comparisons. Every used algebraic claim beyond a
  definition is reconstructed here.
\item
  The controlling book manuscript, Chapter 72, subsection ``Regularity
  and orthogonality on balanced polyhedral spaces,'' particularly
  \texttt{prop:frontier-}\allowbreak{}\texttt{polyhedral-}\allowbreak{}\texttt{bergman-}\allowbreak{}\texttt{positive-}\allowbreak{}\texttt{infinity-}\allowbreak{}\texttt{obstruction}.
  Comparison only, with full displayed proof inspected; no theorem
  premise imported.
\item
  Prior round-05 all-real-reference product argument, its author audit,
  research notes, and embedded evidence/read attestation; and the
  six-cone parent proposal. Fallible research data. The applicable
  finite support and decomposition proofs have been reconstructed in
  Lemmas 1--5. Prior review scope does not certify this new theorem.
\end{enumerate}

\end{document}
